\documentclass{article}
\usepackage[UTF8]{ctex}
\usepackage{amsmath}
\usepackage{amssymb}
\usepackage{graphicx}
\usepackage{geometry}
\geometry{a4paper, margin=2cm}

\newcommand{\problemname}[1]{\begin{center}\Large\bfseries #1\end{center}}
\usepackage{enumitem}

\begin{document}
\problemname{选礼物}

蒜头君要去商店为朋友挑选礼物，他看中了货架上的三个礼物，但他仍然想要拿下来仔细对比一下再决定购买哪一个。由于身高不够高，在旁边有一个高度为 $30$ 的凳子，可以让他站在上面取下货架上的礼品。

\subsection*{输入格式}
输入第一行 $h(100\leh\le150)$ 为蒜头君的身高。

输入第二行 $3$ 个数字 $x_{1},x_{2},x_{3}$，($1\lex_{1},x_{2},x_{3}\le240$)，分别为货架上的三个礼物的高度。

\subsection*{输出格式}
输出一个数字，表示蒜头君能够拿到礼物的个数。

\subsection*{样例输入}
\begin{verbatim}
120
160 142 131
\end{verbatim}

\subsection*{样例输出}
\begin{verbatim}
2
\end{verbatim}


\end{document}
